% Generated draft: requires mathematical and source review.
\documentclass[11pt]{article}
\usepackage[margin=25mm]{geometry}
\usepackage{fontspec}
\setmainfont{Latin Modern Roman}
\usepackage{amsmath,amssymb,mathtools,unicode-math}
\setmathfont{Latin Modern Math}
\usepackage{enumitem,needspace,xurl,hyperref}
\setlength{\parindent}{0pt}
\setlength{\parskip}{5pt}
\setlength{\emergencystretch}{3em}
\newcommand{\N}{\mathbb{N}}
\newcommand{\Z}{\mathbb{Z}}
\newcommand{\Q}{\mathbb{Q}}
\newcommand{\R}{\mathbb{R}}
\newcommand{\C}{\mathbb{C}}
\newcommand{\F}{\mathbb{F}}
\newcommand{\A}{\mathbb{A}}
\newcommand{\PP}{\mathbb{P}}
\newcommand{\E}{\mathbb{E}}
\newcommand{\eps}{\varepsilon}
\newcommand{\dd}{\,\mathrm{d}}
\begin{document}
\section*{UnsolvedMath ID: 20002374}
\textbf{Source problem code:} catalogue entry
\subsection{Definitions and mathematical statement}
A.2 is Andrew Brooke-Taylor's participant-interest statement, not a formal problem. It expresses interest in the implications of continuum-hypothesis work for Woodin's Omega conjecture and in difficulties constructing models with superstrong cardinals, from a then-stated research interest at the \(I3\) level. No exact conjecture, hypotheses, or requested theorem are given. The source's personal research description is historical, not a claim about the participant's present work.

The aid develops a separate forcing diagnostic. Let \(T\) be a theory extending the relevant ZFC background with a countable transitive model \(M\), and assume it is persistent under the set forcing in question: every such generic extension of a \(T\)-model again satisfies \(T\). Then the report gives set-generic extensions with
\[
 M[G]\models T+\mathrm{CH},\qquad
 M[H]\models T+\neg\mathrm{CH},
 \qquad \mathrm{CH}\equiv2^{\aleph_0}=\aleph_1.
\]
A standard CH extension collapses the old continuum to \(\omega_1\) without adding reals, and a Cohen extension adds at least \(\omega_2\) distinct reals while preserving that cardinal. Generics exist externally over the countable transitive model. The required persistence is a hypothesis; it is not asserted for every arbitrary large-cardinal axiom or every fixed cardinal witness.

Consequently generic persistence alone yields neither \(T\models\mathrm{CH}\) nor \(T\models\neg\mathrm{CH}\). Nor can it select a truth value through a consequence relation that is sound for all these generic extensions. This is a soundness diagnostic under the countable-transitive-model assumption, not a new inconsistency theorem or an analysis of every possible Omega-logic axiom.

The aid says Woodin's published conditional anti-CH result requires additional finite Omega-star axiomatizability of the actual theory of \(H(\omega_2)\), the sets of hereditary cardinality below \(\omega_2\). That hypothesis is distinguished from generic stability of a background theory. The exact Omega/Omega-star definitions and full conditional theorem are not supplied, so they are not reconstructed or identified with one another here.

A second contextual distinction is domain-level. \(I3\) is expressed by a nontrivial elementary self-embedding \(j:V_\lambda\to V_\lambda\) for an appropriate rank \(\lambda\). A global superstrong witness has a different domain: an elementary embedding \(J:V\to N\), with critical point \(\kappa\) and \(V_{J(\kappa)}\subseteq N\), in the intended transitive-class formulation. The given rank-domain \(I3\) map is not itself that global map and does not itself constitute an iteration strategy, which supplies coherent branch choices for relevant iteration trees. This does not deny consequences or derived constructions from rank-into-rank hypotheses; none are analyzed by the stated domain observation. No superstrong core model or iteration strategy is constructed.
\subsection{Short English statement}
The source is a participant's research-interest paragraph, not a recoverable conjecture. The aid adds a forcing diagnostic showing why generic persistence alone cannot choose CH, and distinguishes rank-into-rank embeddings from global superstrong witnesses and iteration strategies. It does not resolve Woodin's conjecture or construct a core model.
\subsection{Sources}
\begin{itemize}
\item Local UnsolvedMath registry, permanent ID 20002374, exact source code catalogue entry; canonical \texttt{aim-logic-notes.json}, \texttt{notes[149]}, tag \texttt{section}; AIM Workshop Problem Lists, Logic, \emph{Recent Advances in Core Model Theory}, A.2, Andrew Brooke-Taylor. Created and updated 14 August 2026.
\item Source URL \url{https://aimath.org/WWN/coremodel/coremodel.pdf}. Supplied verification identifies the four-page document, version 29 November 2004, Chapter A ``Participant Contributions'', page 3, between similar personal entries A.1 and A.3. The PDF was not accessed.
\item Woodin's Omega conjecture and a published conditional anti-CH statement are mentioned without a precise publication, theorem number, or full technical hypotheses. No separate authoritative reference is supplied for the forcing diagnostic or embedding comparison.
\item Attempt 1: \texttt{context\_only}, problem status at run \texttt{not\_a\_problem}, rigor label \texttt{proved}, novelty confidence \texttt{low}; classification \texttt{OPEN-TRIAGE}, suggested difficulty L2, registry status \texttt{open}. Candidate: generic-stability CH diagnostic. These are contextual source labels, not a newly formulated open conjecture.
\item Ulam AI / UnsolvedMath Contributors; CC BY 4.0. Machine-generated research aid, not peer reviewed; independently verify before citation.
\end{itemize}
\end{document}
